\documentclass[11pt,a4paper]{article}
% Common preamble for the hint sheets (\input from each file)
\usepackage[utf8]{inputenc}
\usepackage[T1]{fontenc}
\usepackage{lmodern}
\usepackage{amsmath,amssymb,amsthm,mathtools}
\usepackage[a4paper,margin=2.3cm]{geometry}
\usepackage{enumitem}
\usepackage{xcolor}
\usepackage[hidelinks]{hyperref}
\renewcommand{\leq}{\leqslant}
\renewcommand{\le}{\leqslant}
\renewcommand{\geq}{\geqslant}
\renewcommand{\ge}{\geqslant}
\setlength{\parindent}{0pt}
\setlength{\parskip}{0.35em}
% \hintstitle{Sheet title}{Class work / Homework}
\newcommand{\hintstitle}[2]{\begin{center}{\LARGE\bfseries Hints: #1\par}\vspace{0.5em}{\large #2\par}\end{center}
\noindent\textit{The hints to each problem are cumulative. Read one, try again, and open the next only if you are still stuck.}\par\medskip\hrule\medskip}
\newcommand{\problem}[1]{\par\bigskip{\large\bfseries Problem #1}\par\nopagebreak\smallskip}
\newcommand{\subpart}[1]{\par\smallskip\textbf{(#1)}\par\nopagebreak}
\newcommand{\hint}[1]{\par\smallskip\noindent\textcolor{black!70}{\textbf{Hint #1.}}\enspace}
\newcommand{\note}[1]{\par\smallskip\noindent\textit{Note.} #1\par}
\providecommand{\lt}{<}
\providecommand{\gt}{>}
\providecommand{\tightlist}{\setlength{\itemsep}{0pt}\setlength{\parskip}{0pt}}
\emergencystretch=3em

\begin{document}
\hintstitle{Series and infinite products -- 1. Convergence: the standard scale and comparison}{Class work}

\problem{1}
\hint{1} Condensation: cut \([1,\infty)\) into the blocks \(2^k\le n\lt2^{k+1}\). A block has \(2^k\) terms, each between \(a_{2^{k+1}}\) and \(a_{2^k}\), so
\[\tfrac12\sum_{k\ge0}2^{k+1}a_{2^{k+1}}\le\sum_{n\ge1}a_n\le\sum_{k\ge0}2^ka_{2^k}.\]
Integral test: \(f(k+1)\le\int_k^{k+1}f\le f(k)\). The sequence \(D_n=\sum_{k\le n}f(k)-\int_1^nf\) is non-increasing (\(D_{n+1}-D_n=f(n+1)-\int_n^{n+1}f\le0\)) and \(D_n\ge f(n)\ge0\), \(D_1=f(1)\).
\hint{2} For \(a_n=\frac1{n\ln^pn}\) condensation gives \(2^ka_{2^k}=\frac1{(k\ln2)^p}\), so the series converges exactly for \(p\gt1\). For \(\frac1{n\ln n\ln^p\ln n}\) condensation gives terms \(\asymp\frac1{k\ln^pk}\), which is the previous series; again \(p\gt1\). (The integral test does the same with the substitutions \(t=\ln x\), \(t=\ln\ln x\).)
\hint{3} No slowest convergent series. Let \(r_n=\sum_{k\ge n}a_k\) be the tails, \(r_n\downarrow0\), and put \(b_n=a_n/\sqrt{r_n}\). Then \(b_n/a_n\to\infty\), and
\[b_n=\frac{r_n-r_{n+1}}{\sqrt{r_n}}\le2\bigl(\sqrt{r_n}-\sqrt{r_{n+1}}\bigr),\]
which telescopes, so \(\sum b_n\le2\sqrt{r_1}\).
\hint{4} No slowest divergent series. Let \(S_n=\sum_{k\le n}a_k\to\infty\) and \(b_n=a_n/\sqrt{S_n}\), so \(b_n/a_n\to0\). Now \(\sqrt{S_n}-\sqrt{S_{n-1}}=\frac{a_n}{\sqrt{S_n}+\sqrt{S_{n-1}}}\le b_n\), and the left-hand sides add up to \(\sqrt{S_n}-\sqrt{S_1}\to\infty\).

\problem{2}
\hint{1} One direction is a single choice of \(c_n\): which \(c_n\) makes \(\frac{a_n}{c_n}\) and \(\frac{c_n}{b_n}\) equal?
\hint{2} \(c_n=\sqrt{a_nb_n}\) makes both equal to \(\sqrt{a_n/b_n}\), so 2) implies 1).
\hint{3} For 1) \(\Rightarrow\) 2) write \(\sqrt{\frac{a_n}{b_n}}=\sqrt{\frac{a_n}{c_n}\cdot\frac{c_n}{b_n}}\le\frac12\Bigl(\frac{a_n}{c_n}+\frac{c_n}{b_n}\Bigr)\) by AM--GM.

\problem{3}
\hint{1} The series converges. It is enough to show \(a_n\le\frac23n\) for all composite \(n\ge9\); the finitely many terms \(n=4,6,8\) do not matter, and then the terms are at most \((2/3)^n\).
\hint{2} If \(n=ab\) with \(1\lt a\lt b\), then \(a\) and \(b\) are distinct factors of \(b!\), so \(n\mid b!\) and \(a_n\le b\le n/2\). Every composite \(n\) has such a factorisation unless \(n=p^2\) with \(p\) prime (take \(a\) = the least prime divisor).
\hint{3} For \(n=p^2\) the factorial \((2p)!\) contains the factors \(p\) and \(2p\), so \(a_n\le2p=2\sqrt n\), and \(2\sqrt n\le\frac23n\) for \(n\ge9\).

\problem{4}
\hint{1} Compare the term with \(\min(a_n,\frac1n)\): if \(na_n\le1\), then \(\frac{a_n}{1+na_n}\ge\frac{a_n}2\); if \(na_n\ge1\), then \(\frac{a_n}{1+na_n}\ge\frac{a_n}{2na_n}=\frac1{2n}\). So it suffices to show that \(\sum\min(a_n,\frac1n)\) diverges.
\hint{2} \(b_n=\min(a_n,\frac1n)\) is non-increasing, so condensation applies: \(\sum b_n\) converges iff \(\sum_k\min\bigl(2^ka_{2^k},1\bigr)\) does.
\hint{3} If the condensed series converged, its terms would tend to \(0\), so from some point on they equal \(2^ka_{2^k}\), and \(\sum2^ka_{2^k}\) would converge. By condensation again \(\sum a_n\) would converge, a contradiction.

\problem{5}
\hint{1} First check by induction that \(0\lt x_{n+1}\lt x_n\lt1\): \(x_{n+1}=x_n\bigl(1-\frac{x_n}{\sqrt n}\bigr)\) and \(x_n\lt1\le\sqrt n\). Now find the asymptotics of \(x_n\). For \(x'=-x^2/\sqrt t\) one looks at \(1/x\).
\hint{2} \(\frac1{x_{n+1}}-\frac1{x_n}=\frac{x_n-x_{n+1}}{x_nx_{n+1}}=\frac{1}{\sqrt n-x_n}\). Since \(0\lt x_n\lt1\), this is \(\frac1{\sqrt n}\bigl(1+O(n^{-1/2})\bigr)\).
\hint{3} By Stolz, \(\frac{1/x_n}{\sqrt n}\) has the limit \(\lim\frac{1/(\sqrt n-x_n)}{\sqrt{n+1}-\sqrt n}=2\). So \(x_n\sim\frac1{2\sqrt n}\), and \(\sum x_n^\alpha\) converges exactly for \(\alpha\gt2\).

\end{document}
